\section{Divisor Height Laws}
Away from zero, integer heights retain factorization information. These
properties translate classical divisor-function laws \cite{dlmfdivisors} into the
height notation. They concern unweighted filters: changing the positive
weights can preserve contacts while destroying these numerical identities.

\begin{itemize}
\item Divisor factorization and low-height classifications.
\item Square parity and coprime multiplication.
\end{itemize}

\subsection{Factorization and Low Heights}

Let $n\ge2$ have prime factorization $n=\prod_{j=1}^{s}p_j^{e_j}$, with
distinct primes $p_j$ and positive exponents $e_j$. A positive divisor chooses
an exponent independently from $0$ through $e_j$ at each prime. Therefore
\begin{equation*}
\begin{aligned}
h(n)+2&=\tau(n)
&&[\text{Integer Height Identity}]\\
&=\prod_{j=1}^{s}(e_j+1).
&&[\text{Divisor Exponent Choices}]
\end{aligned}
\end{equation*}

This is the divisor-product formula of DLMF equation 27.3.5 \cite{dlmfdivisors}.
Small heights now have exact classifications. The following forms use suitable primes $p$ and $q$:
\begin{equation*}
\begin{aligned}
h(n)=0&\ \Longleftrightarrow\ n=p,\\
h(n)=1&\ \Longleftrightarrow\ n=p^2,\\
h(n)=2&\ \Longleftrightarrow\ n=p^3\text{ or }n=pq\ (p\ne q),\\
h(n)=3&\ \Longleftrightarrow\ n=p^4.
\quad\blacksquare\ \text{[Q.E.D.]}
\end{aligned}
\end{equation*}

Indeed, the possible values of $\tau(n)$ are respectively $2,3,4,5$.
The prime values $2,3,5$ force a single exponent factor; $4$ has the
multiplicative partitions $4$ and $2\cdot2$, giving a cube or two distinct
prime factors. For example, $h(9)=1$, $h(8)=h(10)=2$, and $h(16)=3$.

Low height is thus arithmetic information, not only a visual distance from
the real axis. Height two alone does not distinguish a prime cube from a
product of two distinct primes.


\subsection{Odd Heights and Perfect Squares}

The parity of the height detects a perfect square. Divisors pair as $d$ and
$n/d$; every pair contains two distinct divisors unless $d^2=n$. There is
exactly one unpaired divisor when $n$ is a square, and none otherwise.
Consequently, for every integer $n\ge2$,
\begin{equation*}
\begin{aligned}
\tau(n)\text{ is odd}
&\ \Longleftrightarrow\ n\text{ is a perfect square}
&&[\text{Divisor Pairing}]\\
h(n)\text{ is odd}
&\ \Longleftrightarrow\ n\text{ is a perfect square}.
&&\blacksquare\ [\text{Subtracting Two}]
\end{aligned}
\end{equation*}

This classical divisor-function property is also immediate from the product
formula \cite{dlmfdivisors}: its product is odd exactly when every exponent $e_j$
is even. Thus the parity observation does not assert primality. For example,
$9$ has height one and $16$ has height three; both have already been rejected.


\subsection{Coprime Multiplication}

For coprime integers $a,b\ge2$, the divisor choices are independent. Every
divisor of $ab$ factors uniquely as a divisor of $a$ times a divisor of $b$,
because the two prime factorizations have no common prime. The standard
multiplicativity law \cite{dlmfdivisors}, equation 27.3.1, therefore gives
\begin{equation*}
\begin{aligned}
h(ab)+2&=\tau(ab)
&&[\text{Integer Height Identity}]\\
&=\tau(a)\tau(b)
&&[\gcd(a,b)=1]\\
&=(h(a)+2)(h(b)+2).
&&[\text{Substitution}]
\end{aligned}
\end{equation*}

Equivalently,
\begin{equation*}
\begin{aligned}
h(ab)&=h(a)h(b)+2h(a)+2h(b)+2.
\quad\blacksquare\ \text{[Q.E.D.]}
\end{aligned}
\end{equation*}

The coprimality premise cannot be dropped. For a prime $p$,
$\tau(p^2)=3$, whereas $\tau(p)\tau(p)=4$. Shared prime factors couple the
exponent choices; the height product then requires different bookkeeping.
